\documentclass{article}
\usepackage{amsmath, amssymb}

\begin{document}

\title{CSE400 – Fundamentals of Probability in Computing \\ Lecture L21–L22: Inequalities and Tail Bounds}
\date{}
\maketitle

\section{Introduction and Outline}

The lecture covers:
\begin{itemize}
    \item Tail and motivation of tail bounds
    \item Markov’s Inequality (definition, proof, example)
    \item Chebyshev’s Inequality (definition, proof, example)
    \item Chernoff Bound and Poisson tail
    \item Jensen’s Inequality and convex functions
    \item Case study and system-level understanding
\end{itemize}

\section{Tails of Random Variables}

\subsection{Definition (Tail)}

The tail of a random variable $X$ is defined as:
\begin{equation}
P(X > x)
\end{equation}

This represents the probability that the random variable takes values larger than a threshold $x$.

\subsection{Motivation: Why do we care about tails?}

From the slides:
\begin{itemize}
    \item Jobs delayed more than 24 hours
    \item Packets dropped due to full buffer
\end{itemize}

\subsection{Key Idea}

\begin{itemize}
    \item Mean $\rightarrow$ average behavior (easy to compute)
    \item Tail $\rightarrow$ rare extreme events (hard to compute)
\end{itemize}

Why are tails hard to compute?

\begin{itemize}
    \item Requires summing many small probabilities
\end{itemize}

\section{Tail Probability Examples}

\subsection{Example 1: Load Balancing}

\textbf{Problem Statement:}

Suppose $n$ jobs are distributed among $n$ servers uniformly at random.  
What is the probability that a particular server gets at least $k$ jobs?

\subsection{Modeling}

\begin{itemize}
    \item Jobs assigned randomly
    \item Most servers get few jobs
    \item Occasionally one server gets many jobs
\end{itemize}

This is identified as a tail event (rare overload)

\subsection{Distribution}

\begin{equation}
X \sim \text{Binomial}(n, p)
\end{equation}

\begin{equation}
P(X \geq k) = \sum_{i = k}^{n} \binom{n}{i} p^i (1 - p)^{n - i}
\end{equation}

\subsection{Example 2: Poisson Arrivals}

\textbf{Problem Statement:}

Jobs arrive at a datacenter following a Poisson process with rate $\lambda$.  
Find the probability that at least $k$ jobs arrive in one hour.

\subsection{Modeling}

\begin{itemize}
    \item Jobs arrive randomly over time
    \item Most hours have moderate arrivals
    \item Occasionally bursts occur
\end{itemize}

This is a tail event (sudden spike in load)

\subsection{Distribution}

\begin{equation}
X \sim \text{Poisson}(\lambda)
\end{equation}

\begin{equation}
P(X \geq k) = \sum_{i = k}^{\infty} e^{-\lambda} \frac{\lambda^i}{i!}
\end{equation}

\section{Why Tail Bounds?}

\subsection{Observation}

Exact tail probabilities:
\begin{itemize}
    \item Are hard to compute
    \item Involve long summations
\end{itemize}

\begin{equation}
P(X \geq k) = \sum_{i = k}^{\infty} e^{-\lambda} \frac{\lambda^i}{i!}
\end{equation}

\begin{equation}
P(X \geq k) = \sum_{i = k}^{n} \binom{n}{i} p^i (1 - p)^{n - i}
\end{equation}

\subsection{Key Idea}

Instead of computing exact probabilities:
\begin{itemize}
    \item Compute an upper bound
    \item Easier to compute
    \item Still provides a guarantee
\end{itemize}

\begin{equation}
P(X \geq k) \leq \text{simple bound}
\end{equation}

\textbf{Conclusion:}  
Approximate safely instead of computing exactly

\section{Running Example}

We consider:

\begin{itemize}
    \item Experiment: Flip a fair coin $n$ times
    \item Random variable: $X =$ number of heads
\end{itemize}

\begin{equation}
X \sim \text{Binomial}(n, 1/2)
\end{equation}

\begin{equation}
E[X] = \frac{n}{2}
\end{equation}

\subsection{Tail Event of Interest}

\begin{equation}
P\left(X \geq \frac{3n}{4}\right)
\end{equation}

This corresponds to getting unusually many heads (tail region)

\subsection{Goal}

Develop progressively tighter tail bounds for this probability.

\section{Markov’s Inequality}

\subsection{Key Idea}

\begin{itemize}
    \item Uses only the mean
    \item Large values are rare
\end{itemize}

\subsection{Statement (Markov’s Inequality)}

For a non-negative random variable $X$ and $a > 0$:
\begin{equation}
P(X \geq a) \leq \frac{E[X]}{a}
\end{equation}

\subsection{Interpretation}

\begin{itemize}
    \item Provides a bound using only expectation
    \item Does not require full distribution
\end{itemize}

\subsection{Example (from slide)}

\begin{itemize}
    \item Average = 10
    \item Very large values (e.g., 100+) must be rare
\end{itemize}

(Further slides include detailed proofs, Chebyshev’s inequality, Chernoff bounds, Jensen’s inequality, and examples. Continuing in exact order below.)

\section{Chebyshev’s Inequality}

\subsection{Purpose}

Improves upon Markov’s inequality by incorporating variance.

\subsection{Statement}

For any random variable $X$ with mean $\mu$ and variance $\sigma^2$:
\begin{equation}
P(|X - \mu| \geq k\sigma) \leq \frac{1}{k^2}
\end{equation}

\subsection{Implication}

\begin{itemize}
    \item Controls deviation from the mean
    \item Uses both mean and variance
\end{itemize}

\section{Chernoff Bounds}

\subsection{Purpose}

Provide exponentially tighter bounds for sums of independent random variables.

\subsection{General Idea}

\begin{itemize}
    \item Apply exponential transformation
    \item Use moment generating functions
    \item Optimize bound
\end{itemize}

\subsection{Use Case}

Applied to:
\begin{itemize}
    \item Binomial distributions
    \item Poisson distributions
    \item Sum of independent indicators
\end{itemize}

\subsection{Outcome}

Produces bounds of the form:
\begin{equation}
P(X \geq (1 + \delta)\mu) \leq e^{-(\text{function of } \delta, \mu)}
\end{equation}

\section{Jensen’s Inequality}

\subsection{Convex Function Definition}

A function $f$ is convex if:
\begin{equation}
f(\lambda x + (1 - \lambda)y) \leq \lambda f(x) + (1 - \lambda) f(y)
\end{equation}

\subsection{Jensen’s Inequality}

For a convex function $f$:
\begin{equation}
f(E[X]) \leq E[f(X)]
\end{equation}

\subsection{Use}

\begin{itemize}
    \item Used in proving Chernoff bounds
    \item Relates expectation and transformations
\end{itemize}

\section{Overall Progression of Bounds}

From weakest to strongest:
\begin{itemize}
    \item Markov’s Inequality
    \item Chebyshev’s Inequality
    \item Chernoff Bounds
\end{itemize}

Each step:
\begin{itemize}
    \item Uses more information
    \item Produces tighter bounds
\end{itemize}

\section{Final Key Takeaways}

\begin{itemize}
    \item Tail probabilities represent rare events
    \item Exact computation is difficult due to long summations
    \item Tail bounds provide efficient approximations with guarantees
    \item Different inequalities provide different levels of tightness
    \item These tools are essential in analyzing systems (e.g., networks, queues, load balancing)
\end{itemize}

\end{document}